\begin{trapbox}

	\begin{description}[style=nextline, leftmargin=2.8em, labelsep=0.5em, itemsep=0.35em]
		\item[\textbf{\textcolor{CTrap}{易错点一（忽视真数正负）}}] 运用积或商的对数性质时，未检查 $M$、$N$ 是否为正，或将负数视为真数。
		\item[\textbf{\textcolor{CTrap}{正确理解（真数必须为正）：}}] 对数性质要求所有真数严格大于 $0$。只有 $M,N>0$ 时 $\log_a(MN)=\log_a M+\log_a N$ 才成立。
		\item[\textbf{\textcolor{CTrap}{错因分析（忽视定义域）：}}] 只记公式，无视成立条件，导致在 $M$ 或 $N$ 为负时错误使用性质。
	\end{description}

	\medskip
	\begin{description}[style=nextline, leftmargin=2.8em, labelsep=0.5em, itemsep=0.35em]
		\item[\textbf{\textcolor{CTrap}{易错点二（乱用分配律）}}] 误以为 $\log_a(M+N)=\log_a M+\log_a N$，这是一种常见错误。
		\item[\textbf{\textcolor{CTrap}{正确理解（对数不满足加法分配）：}}] 对数性质只对乘除幂成立，对加法无直接公式。$\log_a(M+N)$ 一般不能化简。
		\item[\textbf{\textcolor{CTrap}{错因分析（受乘法性质干扰）：}}] 将乘法性质的直觉错误迁移到加法上，造成概念混淆。
	\end{description}

	\medskip
	\begin{description}[style=nextline, leftmargin=2.8em, labelsep=0.5em, itemsep=0.35em]
		\item[\textbf{\textcolor{CTrap}{易错点三（换底公式记反）}}] 将换底公式写成 $\log_a b = \frac{\log_c a}{\log_c b}$，分子分母颠倒。
		\item[\textbf{\textcolor{CTrap}{正确理解（分子目标分母底）：}}] 换底公式 $\log_a b = \frac{\log_c b}{\log_c a}$：分子是原真数 $b$ 的对数，分母是原底数 $a$ 的对数。
		\item[\textbf{\textcolor{CTrap}{错因分析（记忆形象不清）：}}] 单纯死记硬背字母位置，缺乏“分子是目标数，分母是底数”的形象化记忆。
	\end{description}

\end{trapbox}
